\section{Positiv}
Grundform (positive degree): \textbf{billig}, \textbf{warm}, \textbf{gut}.

\section{Komparativ}
Vergleich von zwei Dingen (comparative degree): \textbf{billiger}, \textbf{wärmer}, \textbf{besser}. Beispiel: Die Tomaten sind billiger als die Gurken.

\section{Superlativ}
\subsection{Prädikativ und adverbial (Predicative and adverbial)}
Form: \textbf{am + Adjektivstamm + -(e)sten}. Hier wird das Adjektiv nicht dekliniert.
\begin{itemize}
\item Positiv: Die Bohnen sind \textbf{billig}. (The beans are cheap.)
\item Komparativ: Die Tomaten sind \textbf{billiger}. (The tomatoes are cheaper.)
\item Superlativ: Die Gurken sind \textbf{am billigsten}. (The cucumbers are cheapest.)
\item einfach $\rightarrow$ \textbf{am einfachsten}; schnell $\rightarrow$ \textbf{am schnellsten}.
\end{itemize}
\subsection{Attributiv (Attributive)}
Vor einem Nomen: \textbf{-(e)st- + passende Adjektivendung} (depending on article, case, gender, number).
\begin{itemize}
\item Er ist \textbf{der größte Mann}. (He is the tallest man.)
\item Das ist \textbf{das billigste Hotel}. (That is the cheapest hotel.)
\item Sie wohnt in \textbf{dem schönsten Haus}. (She lives in the most beautiful house.)
\end{itemize}
\subsection{Umlaut (Umlaut)}
Einige Adjektive mit a, o, u bekommen einen Umlaut, aber nicht alle (some adjectives take an umlaut, but not all).
\begin{itemize}
\item warm $\rightarrow$ wärmer $\rightarrow$ \textbf{am wärmsten}.
\item klug $\rightarrow$ klüger $\rightarrow$ \textbf{am klügsten}.
\item lang $\rightarrow$ länger $\rightarrow$ \textbf{am längsten}.
\item arm $\rightarrow$ ärmer $\rightarrow$ \textbf{am ärmsten}.
\item groß $\rightarrow$ größer $\rightarrow$ \textbf{am größten}.
\end{itemize}
\subsection{Endung -esten (Extra e)}
Nach s-Lauten (s, ss, ß, z) und vielen Formen auf -t steht häufig oder obligatorisch \textbf{-esten}. Maßgeblich sind auch Aussprache und Wortstruktur (the rule has exceptions).
\begin{itemize}
\item süß $\rightarrow$ \textbf{am süßesten}; kurz $\rightarrow$ \textbf{am kürzesten}.
\item blass $\rightarrow$ \textbf{am blassesten}; hart $\rightarrow$ \textbf{am härtesten}.
\item harmlos $\rightarrow$ \textbf{am harmlosesten}.
\item intelligent $\rightarrow$ \textbf{am intelligentesten}.
\item Ausnahme: groß $\rightarrow$ \textbf{am größten} (nicht: am größesten).
\end{itemize}
\subsection{Unregelmäßige Superlative (Irregular forms)}
\begin{itemize}
\item nah $\rightarrow$ näher $\rightarrow$ \textbf{am nächsten}.
\item gut $\rightarrow$ besser $\rightarrow$ \textbf{am besten}.
\item hoch $\rightarrow$ höher $\rightarrow$ \textbf{am höchsten}.
\item viel $\rightarrow$ mehr $\rightarrow$ \textbf{am meisten}.
\item gern $\rightarrow$ lieber $\rightarrow$ \textbf{am liebsten}.
\end{itemize}

\subsection{Übung 4 (Exercise 4)}
\textbf{Beispiel:} Hotel -- billig: Dieses Hotel ist am billigsten.
\begin{enumerate}
\item Flugzeug -- schnell
\item Koffer -- schwer
\item Sofa -- bequem
\item Film -- langweilig
\item Kleid -- schön
\item Restaurant -- teuer
\item Jacke -- warm
\item Text -- lang
\item Suppe -- scharf
\item Leute -- arm
\item Argument -- dumm
\item Kinder -- klug
\end{enumerate}
\subsection{Übung 5 (Exercise 5)}
\textbf{Beispiel:} Stadtteil -- Wohnungen -- teuer: In diesem Stadtteil sind die Wohnungen am teuersten.
\begin{enumerate}
\item Hotel -- Zimmer -- groß
\item Land -- Berge -- hoch
\item Geschäft -- Personal -- freundlich
\item Museum -- Bilder -- interessant
\item Monat -- Nächte -- kurz
\item Laden -- Gemüse -- frisch
\item Restaurant -- Essen -- lecker
\item Kindergarten -- Kinder -- jung
\item Jahreszeit -- Tage -- heiß
\item Bäckerei -- Brot -- gut
\end{enumerate}
\subsection{Lösungen (Answers)}
\textbf{Übung 4:} a) am schnellsten; b) am schwersten; c) am bequemsten; d) am langweiligsten; e) am schönsten; f) am teuersten; g) am wärmsten; h) am längsten; i) am schärfsten; j) am ärmsten; k) am dümmsten; l) am klügsten.

\textbf{Übung 5:} a) am größten; b) am höchsten; c) am freundlichsten; d) am interessantesten; e) am kürzesten; f) am frischesten; g) am leckersten; h) am jüngsten; i) am heißesten; j) am besten.
